\documentclass{article}

\usepackage{cite} % optional, improves citation formatting

\begin{document}

\title{Bioinformatics Paper Presentation}

\author{Your Name}

\date{}

\maketitle

\section{Introduction}

In this study, large-scale clinico-genomic data was used to analyze mutation-treatment effects in cancer \cite{liu2024_mutation_treatment_effects}.

The dataset and implementation are publicly available \cite{liu2024_precision_cancer2023}.

\section{Conclusion}

This demonstrates how computational methods can improve precision medicine.

\bibliographystyle{plain}

\bibliography{references} % <-- this links to references.bib

\documentclass{article}

\usepackage{cite} % optional, improves citation formatting

\begin{document}

\title{Bioinformatics Paper Presentation}

\author{Your Name}

\date{}

\maketitle

\section{Introduction}

In this study, large-scale clinico-genomic data was used to analyze mutation-treatment effects in cancer \cite{liu2024_mutation_treatment_effects}.

The dataset and implementation are publicly available \cite{liu2024_precision_cancer2023}.

I thought this paper would be a great introduction to implimenting Bioinformatics papers. Not too complext and offered room to expand on the material.

\section{Conclusion}

This demonstrates how computational methods can improve precision medicine.

\bibliographystyle{plain}

\bibliography{references} 


\documentclass[12pt]{article}

\usepackage[margin=1in]{geometry}
\usepackage{setspace}
\usepackage{cite}
\usepackage{hyperref}

\title{Analysis of an Open-Source Framework for End-to-End Electronic Health Record Data}
\date{}

\begin{document}

\maketitle

\onehalfspacing

\section{Introduction}

Electronic health record (EHR) data has become an important resource for modern biomedical research. 
However, analyzing this data is challenging due to its complexity, scale, and variability across healthcare systems. 
Recent work has focused on developing computational frameworks that streamline the process of extracting, processing, 
and analyzing EHR data.

In this paper, the authors present an open-source framework designed to enable end-to-end analysis of EHR data

This data shows a useful tool and analysis methods on EHR data. however, due to blocked and private datasets, implimentation was not possible for 
recreating the whole paper. Only supplemental figure 10 was possible for the paper recreation.


\cite{heumos2024open}. 

\documentclass[12pt]{article}

\usepackage{cite}
\usepackage{hyperref}

\title{Analysis of Tumor-Normal Single-Cell Ecosystems Across Cancer Types}
\date{}

\begin{document}

\maketitle

\onehalfspacing

\section{Introduction}

Understanding tumor biology at the single-cell level has become a central focus in cancer research. 
Traditional bulk sequencing methods average signals across many cells, which can obscure important cellular diversity.

In this study, the authors perform a large-scale single-cell analysis across nearly one thousand tumors spanning 30 cancer types.
This paper was harder to impliment independently, but was still very informative on sequencing methods and computational biology.


\cite{kang2024systematic}. Their goal is to better understand tumor microenvironments and how cancer cells interact with surrounding normal and immune cells.



\bibliographystyle{plain}
\bibliography{references}


\documentclass[12pt]{article}

\usepackage{cite}
\usepackage{hyperref}

\title{Analysis of Tumor-Normal Single-Cell Ecosystems Across Cancer Types}
\date{}

\begin{document}

\maketitle

\onehalfspacing

\section{Introduction}

Understanding tumor biology at the single-cell level has become a central focus in cancer research. 
Traditional bulk sequencing methods average signals across many cells, which can obscure important cellular diversity.

In this study, the authors perform a large-scale single-cell analysis across nearly one thousand tumors spanning 30 cancer types.
This paper was harder to impliment independently, but was still very informative on sequencing methods and computational biology.


\cite{kang2024systematic}. Their goal is to better understand tumor microenvironments and how cancer cells interact with surrounding normal and immune cells.



\bibliographystyle{plain}
\bibliography{references}

\end{document}
\end{document}
